% Historical source (removed by author request): references/_OLD_Übungsaufgaben/Übungsblatt 3 - Vektorechnung/Blatt3.tex:47
\begin{ExerciseBlock}
\ExerciseHeading{03}{Skalarprodukt und Vektorprodukt}{vector_algebra_scalar_and_cross_products}
Gegeben seien die folgenden Vektoren 
$\vec{a}=\left(\begin{array}{c} 3 \\ 2 \\ -4 \end{array}\right)$, 
$\vec{b}=\left(\begin{array}{c} -2 \\ 0 \\ 4\end{array}\right)$ und
${\vec{c}}=\left(\begin{array}{c} -5 \\ 1 \\ 4 \end{array}\right)$. \par

Berechnen Sie
\begin{enumerate}[label=\alph*)]
	\item $\vec{a} \cdot \vec{b} = $
	\item $\vec{a} \times \vec{b} = $
	\item $\vec{b} \times \vec{a} = $
	\item $\vec{a} \cdot \left( \vec{b} \times \vec{c} \right) =$
	\item $\vec{b} \cdot \left( \vec{c} \times \vec{a} \right) =$
	\item $\vec{a} \times \left( \vec{b} \times \vec{c} \right) =$
	\item $\vec{b} \left( \vec{a} \cdot \vec{c} \right) - \vec{c} \left( \vec{a} \cdot \vec{b} \right) = $
\end{enumerate} 

\end{ExerciseBlock}
\UseExerciseSolution{vector_algebra_scalar_and_cross_products}
